
% ============================================================================
\clearpage
\section*{v4.8.0-r15 Source-Only: The Octonionic Hopf Fibration as Carrier of the Triality Slots}
\addcontentsline{toc}{section}{v4.8.0-r15 Source-Only: Octonionic Hopf Fibration as Carrier of the Triality Slots}
\label{sec:v480-r15-hopf-triality-carrier}
% ============================================================================

\paragraph{Purpose.}
Layer r14 proposed the three triality slots of \(\operatorname{Re}((xy)z)\) as generation index but left their
carrier open.  This layer shows that the octonionic Hopf fibration of Theorem~\ref{thm:hopf}, which the Companion
already attaches to \(S_2\cong\OO P^1\) and to three generations, carries the slots.  Model construction only; no
prediction is made and no data are compared.  Earlier layers are not rewritten.

\noindent\textbf{Source-only status.}
This v4.8.0-r15 layer is source-only.  It does not build a PDF, does not assign a new DOI,
and does not authorize a public release.  The v4.8.0 Version DOI
\href{https://doi.org/10.5281/zenodo.23365244}{10.5281/zenodo.23365244} and the Concept DOI
\href{https://doi.org/10.5281/zenodo.19210728}{10.5281/zenodo.19210728} are unchanged by it.

\subsection*{v4.8.0-r15 unit triples and the Hopf map}
\label{subsec:v480-r15-triples}

\paragraph{Theoremlet v4.8.0-r15.1.}
Let \(\mathcal T=\{(a,b,c)\in(S^7)^3:(ab)c=1\}\).  Then \(\mathcal T\) is a fourteen-dimensional manifold
(\(a,b\) free, \(c=\overline{ab}\)), closed under the cyclic shift \((a,b,c)\mapsto(b,c,a)\) and under
\((a,b,c)\mapsto(\bar a,\bar c,\bar b)\), and invariant under \(G_2\) acting diagonally.  For the Hopf map
\(h(x,y)=(|x|^2-|y|^2,\,2x\bar y)\colon S^{15}\to S^8\), the points over the equator \(|x|=|y|=1/\sqrt2\) give
unit triples \((\sqrt2\,x,\sqrt2\,\bar y,\bar v)\in\mathcal T\) with \(v=2x\bar y\), and the fibre over \(v\) is
\(\{(a,\bar a v)\}\cong S^7\).

\paragraph{Proof sketch.}
\(c=\overline{ab}\) gives \((ab)c=|ab|^2=1\).  By alternativity \(b(\bar b\bar a)=\bar a\), so \(bc=\bar a\) and
\((bc)a=1\); the remaining cyclic and reflected identities follow in the same way.  Elements of \(G_2\) preserve
products and conjugation.  The Hopf statements are the norm identity \((|x|^2-|y|^2)^2+4|x|^2|y|^2=1\) and
\(a(\bar a v)=v\).  All identities were also checked numerically (script \texttt{k29\_hopf\_triality\_traeger.py}).
\(\square\)

\paragraph{Theoremlet v4.8.0-r15.2 (symmetric point).}
\((a,a,a)\in\mathcal T\) if and only if \(a^3=1\), i.e.\ \(a=\exp(2\pi u/3)\) for an imaginary unit \(u\).  For
\(u=e_7\) the point is \(\omega\) in \(\mathrm{span}(1,e_7)\), the complex plane at the end of the cascade, and
its stabilizer in \(G_2\) is the colour group \(SU(3)\) (dimension \(8\)).

\paragraph{Proof sketch.}
\(a\) and its powers lie in the commutative subalgebra \(\mathrm{span}(1,u)\cong\mathbb C\); the stabilizer of
\(a\) in \(G_2\) is the stabilizer of \(u\).  The dimension was computed from the derivation algebra.  \(\square\)

\subsection*{v4.8.0-r15 reading and a check}
\label{subsec:v480-r15-reading}

The two components of the Hopf total space and the base point are the three triality slots; the triole is their
cyclic symmetry, and \(G_2\) acts as the colour side.  Choosing the symmetric generation point fixes \(u\) and breaks
\(G_2\) to the colour \(SU(3)\).  This gives a natural carrier for the assignment of r14; it does not derive the
existence of three generations.

Check: in \(\mathrm{span}(1,e_7)\) the condition \((ab)c=1\) becomes \(\alpha_1+\alpha_2+\alpha_3\equiv0\pmod{2\pi}\)
for the phases.  Reading the Koide phases \(\delta+2\pi k/3\) of layer r10 as these phases requires
\(3\delta\equiv0\), against \(\delta=0.2222\) for the charged leptons.  The direct reading fails.  That
\(3\delta\approx2/3\) is the known value \(2/9\cdot3\), whose generality was refuted in r11; it is not counted.  How
a mass hierarchy and \(K=2/3\) could arise from deviations from the symmetric triple remains open.

\subsection*{v4.8.0-r15 ledger and guard update}
\label{subsec:v480-r15-ledger-guard}

\[
R_{\rm cg}^{v480\text{-}r15}=(0,0,1)
\]
for (separation consistent with \(G_2\); carrier of the triality slots identified; mass mechanism open).

The layer spends no guard:
\[
R_{\rm glob}^{(v480\text{-}r15,+)}=(0,1,1,1,1),
\qquad
\Gamma_{\rm guard}^{v480\text{-}r15}=(1,1,1,1,1),
\qquad
\beta_4=1/256\ \text{unspent}.
\]
It does not imply exact finite \(C5\) zero closure, exact all-readout terminality,
\(C6/T2\), a global photon-clock/RH theorem, or global all-readout closure.
